%%%%%%%%%%%%%%%%
% PROBLEMA
%%%%%%%%%%%%%%%%

\section*{Resumen}

El proyecto aborda la necesidad de reducir el esfuerzo repetitivo y las diferencias estructurales entre proyectos que enfrentan los equipos de desarrollo al iniciar nuevos proyectos de software e integrar funcionalidades compartidas en entornos bancarios como el Laboratorio digital del Banco W. Se propone el diseño, implementación y validación de una herramienta CLI para entornos Node.js que automatice el arranque de proyectos backend, estandarice su estructura desde las fases tempranas e integre capacidades reutilizables de forma desacoplada. La herramienta contempla funcionalidades como la generación de proyectos a partir de plantillas disponible para los equipo, un mecanismo de composición que permite incorporar módulos compartidos siguiendo reglas definidas y la generación automática de documentación OpenAPI. Su diseño se fundamenta en principios de arquitectura hexagonal y Domain-Driven Design, y se alinea con estándares del sector financiero como BIAN e ISO 20022. 
\\\\
\textbf{Palabras clave:} CLI, Reutilización de software, Scaffolding de proyectos, Arquitectura hexagonal, Domain-Driven Design, BIAN

\pagebreak

\section*{Abstract}

This project addresses the need to reduce repetitive effort and structural differences between projects faced by development teams when starting new software projects and integrating shared functionalities in banking environments such as the Digital Laboratory of Banco W. The design, implementation, and validation of a CLI tool for Node.js environments is proposed, aiming to automate backend project bootstrapping, standardize project structure from early stages, and integrate reusable capabilities in a decoupled manner. The tool includes features such as project generation from templates available to development teams, a composition mechanism that allows incorporating shared modules following defined rules, and automatic OpenAPI documentation generation. Its design is grounded in hexagonal architecture and Domain-Driven Design principles, and aligns with financial sector standards such as BIAN and ISO 20022.
\\\\
\textbf{Keywords:} CLI, Software reuse, Project scaffolding, Hexagonal architecture, Domain-Driven Design, BIAN